\chapter{Introduction}\label{chap:Intro}

\section{L'entreprise Obeo}
\obeo{} est une société de service et un éditeur de logiociels. Son expertise dans le domaine de l'ingéniérie des modèles (démarche MDA), lui permet de proposer des solutions allant de la création à la refonte d'applications informatiques. 


\begin{figure}[htb]
  \centering
  \includegraphics[scale=.4]{img/logoobeo.eps}
  \label{fig:obeo}
\end{figure}

Ces solutions proposées permettent notamment de diminuer les délais de projets et de diminuer les risques d'erreurs. Les amélioration des performances d'adaption et d'agilité font aussi parties des objectifs des outils et méthode de la société \obeo{}.




La société \obeo{} est aussi un membre actif dans le domaine Open Source et est membre de la fondation Eclipse. Elle est à l'initiative du projet Acceleo, un générateur de code basé sur le framework EMF. 


\section{Objectifs}
Dans le cadre du module \emph{SC Project}, nous avons eu l'opportunité de travailler sur la conception d'outils de modélisation dédiés aux applications web. En partant de trois méta models mis à disposition par \textsc{Obeo}
(\cf{} section \ref{sec:met}) l'objectif était de créer des modèles et les générateurs de code associé pour produire le code d'une aplication
% LocalWords:  TP First Rom
